\section{Rosser's strengthening}
\label{sec:first-incompleteness-04}

\begin{theorem}[Rosser]
\label{thm:rosser}
Let \(T\) be a consistent computably enumerable theory extending \(\mathsf Q\).
Then there is
a sentence \(R_T\) such that
\[
  T\notderives R_T
  \qquad\text{and}\qquad
  T\notderives\neg R_T.
\]
Thus consistency alone suffices for deductive incompleteness.
\textit{Status: complete proof.}  The two full finite-comparison
branches are in \cref{annex:p5.inc.rosser-chain}, using the selected
certified interface of \cref{thm:p05-rep-certified-proof}.
\end{theorem}

The original source is \cite{Rosser1936}.  The proof here is the
explicit modern \(\mathsf Q\)-based reconstruction with the local
predicate selection just specified, not an identification with the
historical paper's system or proof numbering.

\begin{proofplan}
By \cref{thm:diagonal-lemma}, fix a sentence \(R_T\) with
\[
  \mathsf Q\derives R_T\leftrightarrow\forall y\,\bigl(
  \mathsf{Proof}_T(y,\overline{\ulcorner R_T\urcorner})\to
  \exists z\leq y\,
  \mathsf{Proof}_T(z,\overline{\ulcorner\neg R_T\urcorner})\bigr),
\]
so that \(R_T\) says: every proof of \(R_T\) has a proof of \(\neg R_T\)
with no larger code.  The bounded comparison uses the registered
object-language abbreviation \(\leq\).

First branch: suppose \(T\derives R_T\), and let \(p\) be the least code
of such a proof.  Numeralwise representation of the proof relation gives
\(T\derives\mathsf{Proof}_T(\bar p,\overline{\ulcorner R_T\urcorner})\),
so the fixed point yields
\(T\derives\exists z\leq\bar p\,
\mathsf{Proof}_T(z,\overline{\ulcorner\neg R_T\urcorner})\).
Consistency gives \(T\notderives\neg R_T\), so numeralwise representation
gives
\(T\derives\neg\mathsf{Proof}_T(\bar q,
\overline{\ulcorner\neg R_T\urcorner})\) for each \(q\leq p\).
\Cref{lem:q-bounded-cases} turns the bounded existential into the finite
disjunction over \(z=\bar0,\ldots,\bar p\), each disjunct refuted by the
preceding line, so \(T\) is inconsistent, contradicting consistency.  No
metatheoretic proof of \(\neg R_T\) is produced at any point.

Second branch: suppose \(T\derives\neg R_T\), and let \(q\) be the least
code of such a proof; numeralwise representation gives
\(T\derives\mathsf{Proof}_T(\bar q,
\overline{\ulcorner\neg R_T\urcorner})\).  Reason inside \(T\) from the
assumption \(\mathsf{Proof}_T(y,\overline{\ulcorner R_T\urcorner})\).
\Cref{lem:q-dichotomy} gives \(y\leq\bar q\lor\bar q\leq y\).  In the
first case, \cref{lem:q-bounded-cases} gives
\(y=\bar0\lor\dots\lor y=\bar q\); consistency gives
\(T\notderives R_T\), so numeralwise representation gives
\(T\derives\neg\mathsf{Proof}_T(\bar k,
\overline{\ulcorner R_T\urcorner})\) for each \(k\leq q\); every
disjunct is contradicted, and the bounded existential follows from the
contradiction.  In the second case, \(\bar q\leq y\) makes
\(z\defeq\bar q\) a witness for
\(\exists z\leq y\,
\mathsf{Proof}_T(z,\overline{\ulcorner\neg R_T\urcorner})\).
Discharging the assumption and applying the fixed point yields
\(T\derives R_T\), so \(T\) is inconsistent, again contradicting
consistency.

The representation and diagonal dependencies are supplied by
\cref{thm:p05-rep-certified-proof,thm:p05-inc-diagonal}; the complete
two-branch derivation is in \cref{annex:p5.inc.rosser-chain}.
The bounded-comparison steps themselves are carried by
\cref{lem:q-bounded-cases,lem:q-dichotomy}, which have complete proofs.
\end{proofplan}

\begin{corollary}[Undecidability of theoremhood]
If \(T\) is a consistent computably enumerable extension of \(\mathsf Q\),
then the set of codes of theorems of \(T\) is computably enumerable but
undecidable.
\textit{Status: complete proof.}  Finite-extension consistency and
the complete effective-extension argument are in
\cref{annex:p5.inc.finite-extension,annex:p5.inc.theoremhood-chain}.
\end{corollary}

\begin{proofplan}
The theorem set is computably enumerable by enumeration of finite \(T\)-proofs.  If it were
decidable, then for each sentence \(\sigma\) one could decide whether
\(T\derives\sigma\).  Run this decider on both \(R_T\) and \(\neg R_T\); the
two negative answers are compatible with decidability, so this alone is not a
contradiction.  Instead enumerate the sentences as
\(\sigma_0,\sigma_1,\ldots\).  Starting with the empty finite set \(D_0\), let
\(D_{n+1}=D_n\cup\{\sigma_n\}\) if \(T\cup D_n\cup\{\sigma_n\}\) is
consistent, and otherwise let
\(D_{n+1}=D_n\cup\{\neg\sigma_n\}\).  The consistency test is decidable:
by the deduction theorem, inconsistency of the finite extension is equivalent
to \(T\) proving the negation of the conjunction of its finitely many added
sentences.  At least one choice preserves consistency, since otherwise
\(T\cup D_n\) would prove both \(\sigma_n\) and \(\neg\sigma_n\).

The theory \(T^*=T\cup\bigcup_nD_n\) is computably enumerable, consistent, extends
\(\mathsf Q\), and is deductively complete because a side of every sentence
is among its axioms.  Rosser's theorem applied to \(T^*\) gives an independent
sentence, a contradiction.
The complete finite-extension proof in
\cref{annex:p5.inc.finite-extension} checks the actual nonvacuous
discharge and quantifier conditions, and
\cref{annex:p5.inc.theoremhood-chain} supplies the effective construction
and its contradiction with the fully proved Rosser theorem.
\end{proofplan}

\begin{remark}
An alternative proof reduces the self-application domain set to theoremhood by uniform
representability.  The proof above makes explicit why independence of one
fixed sentence does not by itself immediately contradict decidability of a
set.
\end{remark}
